\documentclass[UTF8,11pt]{ctexart}
\usepackage[a4paper,margin=22mm]{geometry}
\usepackage{graphicx,enumitem,amssymb,tabularx}
\usepackage[normalem]{ulem}
\usepackage{xeCJKfntef}
\setlength{\parindent}{0pt}
\setlength{\parskip}{.6em}
\sloppy
\begin{document}
\section*{2023年6月英语四级真题(第1套)}
2023 年6 月大学英语四级考试真题（第1 套）


\subsection*{\textbf{Part I Writing} \textbf{( 30 minutes)}}


\textbf{Directions:} \textit{Suppose the student union of your university is organizing an online discussion on interpersonal} \textit{relationships. You are to write an essay on ways to maintain a warm and friendly relationship with your} \textit{classmates and on the benefits of such a relationship. You will have 30 minutes for the task. You should write} \textit{at least }\textit{\uline{120 }}\textit{words but no more than }\textit{\uline{180 }}\textit{words.}


\subsection*{\textbf{Part II} \textbf{Listening Comprehension} \textbf{( 25 minutes)}}


\subsection*{\textbf{Section A}}


\textbf{Directions:} \textit{In this section, you will hear three news reports. At the end of each news report, you will hear two} \textit{or three questions. Both the news report and the questions will be spoken only once. After you hear a question,} \textit{you} \textit{must} \textit{choose} \textit{the} \textit{best} \textit{answer} \textit{from} \textit{the} \textit{four} \textit{choices} \textit{marked} \textit{A),} \textit{B),} \textit{C)} \textit{and} \textit{D).} \textit{Then} \textit{mark} \textit{the} \textit{corresponding letter on }\textbf{\textit{Answer Sheet 1}}\textit{ with a single line through the centre.}


\textbf{Questions 1 and 2 are based on the news report you have just heard.}


1.


\begin{itemize}[leftmargin=2em,itemsep=.2em,topsep=.4em]\item[A.] A man was taken to a hospital.
\item[B.] A man was bitten by a snake.
\item[C.] A man fell off his toilet seat.
\item[D.] A man kept a 4-foot snake as a pet.
\end{itemize}

2.


\begin{itemize}[leftmargin=2em,itemsep=.2em,topsep=.4em]\item[A.] Who owned the snake.
\item[B.] How the snake was captured.
\item[C.] Whether the snake was infected.
\item[D.] Where the snake had been taken.
\end{itemize}

\textbf{Questions 3 and 4 are based on the news report you have just heard.}


3.


\begin{itemize}[leftmargin=2em,itemsep=.2em,topsep=.4em]\item[A.] Posting her daughter’s photos on social media.
\item[B.] Sharing her photos with famous movie stars.
\item[C.] Amusing herself by going to ball parties.
\item[D.] Taking her trash out in fancy dresses.
\end{itemize}

\noindent\begin{tabularx}{\linewidth}{@{}lXXXX@{}}4. & \textbf{A.} To record her achievements. & \textbf{B.} To please her daughter. & \textbf{C.} To amuse people. & \textbf{D.} To make herself popular.\\\end{tabularx}\par

\textbf{Questions 5 to 7 are based on the news report you have just heard.}


5.


\begin{itemize}[leftmargin=2em,itemsep=.2em,topsep=.4em]\item[A.] Get a free meal after answering some questions.
\item[B.] Have a meal even if they have no money.
\item[C.] Have a chance of winning a \$100 prize.
\item[D.] Eat as much as they want for \$10.
\end{itemize}

6.


\begin{itemize}[leftmargin=2em,itemsep=.2em,topsep=.4em]\item[A.] It was brought up by two staffers.
\item[B.] It was suggested by some of her customers.
\item[C.] It originated from a donation to her staff.
\item[D.] It helped to popularize her restaurant.
\end{itemize}

7.


\begin{itemize}[leftmargin=2em,itemsep=.2em,topsep=.4em]\item[A.] Fifty customers have offered donations.
\item[B.] Most staffers have received messages of kindness.
\item[C.] Many people have come to eat at the restaurant.
\item[D.] More people have been giving than taking.
\end{itemize}

\subsection*{\textbf{Section B}}


\textbf{Directions:} \textit{In this section, you will hear two long conversations. At the end of each conversation, you will} \textit{hear four questions. Both the conversation and the questions will be spoken only once. After you hear a} \textit{question, you must choose the best answer from the four choices marked A), B), C) and D). Then mark the} \textit{corresponding letter on }\textbf{\textit{Answer Sheet 1}}\textit{ with a single line through the centre.}


\textbf{Questions 8 to 11 are based on the conversation you have just heard.}


8.


\begin{itemize}[leftmargin=2em,itemsep=.2em,topsep=.4em]\item[A.] He is a famous writer.
\item[B.] He is a psychologist.
\item[C.] He is a host for a TV program.
\item[D.] He is a primary school teacher.
\end{itemize}

9.


\begin{itemize}[leftmargin=2em,itemsep=.2em,topsep=.4em]\item[A.] Why friends break off contact all of a sudden.
\item[B.] Why people fail to respond to emails promptly.
\item[C.] Why parents raise their children in different ways.
\item[D.] Why social media accounts vanish without a trace.
\end{itemize}

10.


\begin{itemize}[leftmargin=2em,itemsep=.2em,topsep=.4em]\item[A.] They attempt to ignore the whole situation.
\item[B.] They scream to get their parents back.
\item[C.] They avoid showing their emotions.
\item[D.] They simply shut themselves down.
\end{itemize}

11.


\begin{itemize}[leftmargin=2em,itemsep=.2em,topsep=.4em]\item[A.] They try to express their feelings and thoughts effectively.
\item[B.] They are on better terms with friends and romantic partners.
\item[C.] They may regard any difference as the end of a relationship.
\item[D.] They attach more value to their relationships with others.
\end{itemize}

\textbf{Questions 12 to 15 are based on the conversation you have just heard.}


\noindent\begin{tabularx}{\linewidth}{@{}lXXXX@{}}12. & \textbf{A.} Their quality. & \textbf{B.} Their design. & \textbf{C.} Their price. & \textbf{D.} Their color.\\\end{tabularx}\par

13.


\begin{itemize}[leftmargin=2em,itemsep=.2em,topsep=.4em]\item[A.] Jeans are available in a greater variety in America.
\item[B.] America has the best weaving tools in the world.
\item[C.] America makes the best-known brands of jeans.
\item[D.] Jeans are a typical American garment.
\end{itemize}

\noindent\begin{tabularx}{\linewidth}{@{}lXXXX@{}}14. & \textbf{A.} They are artificial. & \textbf{B.} They are unique. & \textbf{C.} They are natural. & \textbf{D.} They are special.\\\end{tabularx}\par

15.


\begin{itemize}[leftmargin=2em,itemsep=.2em,topsep=.4em]\item[A.] They are worth the price.
\item[B.] They are much too pricey.
\item[C.] They are popular with boys.
\item[D.] They are for casual wearing.
\end{itemize}

\subsection*{\textbf{Section C}}


\textbf{Directions:}\textit{ In this section, you will hear three passages. At the end of each passage, you will hear three or} \textit{four questions. Both the passage and the questions will be spoken only once. After you hear a question, you} \textit{must choose the best answer from the four choices marked A), B), C) and D). Then mark the corresponding} \textit{letter on }\textbf{\textit{Answer Sheet 1}}\textit{ with a single line through the centre.}


\textbf{Questions 16 to 18 are based on the passage you have just heard.}


16.


\begin{itemize}[leftmargin=2em,itemsep=.2em,topsep=.4em]\item[A.] He wants to see it again.
\item[B.] He desires more in life.
\end{itemize}

\begin{itemize}[leftmargin=2em,itemsep=.2em,topsep=.4em]\item[C.] He longs to become a superstar himself.
\item[D.] He feels as inspired as other audience members.
\end{itemize}

17.


\begin{itemize}[leftmargin=2em,itemsep=.2em,topsep=.4em]\item[A.] It is extremely artistic.
\item[B.] It is rather unrealistic.
\item[C.] It is somewhat complicated.
\item[D.] It is relatively predictable.
\end{itemize}

18.


\begin{itemize}[leftmargin=2em,itemsep=.2em,topsep=.4em]\item[A.] They are too simple.
\item[B.] They are full of shootings.
\item[C.] They are basically misleading.
\item[D.] They are biased against women.
\end{itemize}

\textbf{Questions 19 to 21 are based on the passage you have just heard.}


19.


\begin{itemize}[leftmargin=2em,itemsep=.2em,topsep=.4em]\item[A.] It is an intuitive way to solidify friendship.
\item[B.] It is a means to inspire creative thinking.
\item[C.] It can help connect people.
\item[D.] It can highlight leadership.
\end{itemize}

20.


\begin{itemize}[leftmargin=2em,itemsep=.2em,topsep=.4em]\item[A.] Make them more open to learning.
\item[B.] Allow them to recite data points.
\item[C.] Enable them to remember the main idea.
\item[D.] Stimulate them to engage in discussions.
\end{itemize}

21.


\begin{itemize}[leftmargin=2em,itemsep=.2em,topsep=.4em]\item[A.] Enrich their own experience.
\item[B.] Inspire listeners’ imagination.
\item[C.] Explain insightful ideas.
\item[D.] Convey fundamental values.
\end{itemize}

\textbf{Questions 22 to 25 are based on the passage you have just heard.}


22.


\begin{itemize}[leftmargin=2em,itemsep=.2em,topsep=.4em]\item[A.] The number of immigrants to the U.S. is declining.
\item[B.] Another wave of immigrants is hitting the U.S.
\item[C.] Immigrants have been contributing to the U.S.
\item[D.] Immigrants outnumber U.S.-born Americans.
\end{itemize}

23.


\begin{itemize}[leftmargin=2em,itemsep=.2em,topsep=.4em]\item[A.] They have founded most Fortune 500 companies.
\item[B.] They have fewer chances to be hired by U.S. companies.
\item[C.] More of them expect their children to succeed in business.
\item[D.] More of them are successful business people.
\end{itemize}

24.


\begin{itemize}[leftmargin=2em,itemsep=.2em,topsep=.4em]\item[A.] Thirty-four percent of them use credit for their daily purchases.
\item[B.] Their level of debt is lower than that of native-born Americans.
\item[C.] Nineteen percent of them borrow money from friends and family.
\item[D.] They have higher installment loan debt than native-born Americans.
\end{itemize}

25.


\begin{itemize}[leftmargin=2em,itemsep=.2em,topsep=.4em]\item[A.] Collaborate with native-born Americans.
\item[B.] Borrow money from financial institutions.
\item[C.] Keep their traditional values and old habits.
\item[D.] Find employment in competitive businesses.
\end{itemize}

\subsection*{\textbf{Part III Reading Comprehension ( 40 minutes)}}


\subsection*{\textbf{Section A}}


\textbf{Directions:} \textit{In this section, there is a passage with ten blanks. You are required to select one word for each} \textit{blank from a list of choices given in a word bank following the passage. Read the passage through carefully} \textit{before making your choices. Each choice in the bank is identified by a letter. Please mark the corresponding} \textit{letter for each item on }\textbf{\textit{Answer Sheet 2}}\textit{ with a single line through the centre. You may not use any of the words} \textit{in the bank more than once.}


You probably haven’t taken the time to think of all the work that went into creating the shirt on your back. I mean, how hard it could be to create fabric and \uline{26} it into a shirt shape. Don’t machines do all that? Well, creating fabric from cotton, which is the most \uline{27} clothing material, is actually a process that involves a lot of water, 2,700 liters per shirt to be \uline{28 }. Take a look at the video below from \textit{National Geographic} for some more mind-blowing \uline{29} about cotton clothing production.


Clean water is \uline{30} becoming one of the most sought-after resources in the world. Given how large the \uline{31} and cotton industries are, they take up a lot of our fresh water demands across the world, according to \textit{The Huffington Post}. The video from \textit{National Geographic} was created to spread \uline{32} of how environmentally harmful cotton is. But the situation can be made better. Through better water management and farming practices, water usage in cotton production can be cut down by \uline{33} 40 percent.


Called “Better Cotton”, this environmentally conscious product will save millions of liters of water a year simply from \uline{34} the demands of cotton production. Cotton doesn’t have to go, since it is, after all, one of the most useful cash crops across the globe. However, as water supplies \uline{35 }, farmers and consumers need to be more conscious of the effect that these products have on the environment as a whole.


\begin{itemize}[leftmargin=2em,itemsep=.2em,topsep=.4em]\item[A.] abstracts
\item[B.] abundant
\item[C.] awareness
\item[D.] conscience
\item[E.] exact
\item[F.] increasingly
\item[G.] intense
\item[H.] mend
\item[I.] nearly
\item[J.] reckoning
\item[K.] reducing
\item[L.] sew
\item[M.] shrink
\item[N.] statistics
\item[O.] textile
\end{itemize}

\subsection*{\textbf{Section B}}


\textbf{Directions:} \textit{In this section, you are going to read a passage with ten statements attached to it. Each statement} \textit{contains information given in one of the paragraphs. Identify the paragraph from which the information is} \textit{derived. You may choose a paragraph more than once. Each paragraph is marked with a letter. Answer the} \textit{questions by marking the corresponding letter on }\textbf{\textit{Answer Sheet 2}}\textit{.}


\subsection*{\textbf{The Spoken Web}}


A) We’re growing more used to chatting to our computers, phones and smart speakers through voice assistants like Amazon’s Alexa, Apple’s Siri and Microsoft’s Cortana. Blind and partially sighted people have been using text-to-speech converters for decades.


B) Out of these assistants, Siri is the most well-known. The assistant uses voice inquiries and a natural-language user interface (界面) to answer questions. The software adapts to users’ individual language usages, searches, and preferences, with continuing use.


C) Some think voice could soon take over from typing and clicking as the main way to interact online.


But what are the challenges of moving to “the spoken web”?


D) What use is written online content if you can’t read? That is the situation facing illiterate (不识字) African farmers. They are often denied crucial information the web offers many others. With a literacy rate in some parts of Africa at only 22.6\%, farmers are often “underpaid for their produce because they might be unaware of the prevailing prices,” says Francis Dittoh, a researcher behind Mr Meteo, a speech-based weather information service.


E) “The most frequently heard complaint is about rainfall predictions,” says Mr Dittoh, who lives in Tamale, northern Ghana. “They tell us the methods their forefathers used to predict the weather don’t seem to work as well these days.” This is down to climate change, he believes. Yet knowing when it’s going to rain is vital for farmers wanting to sow seeds, irrigate crops or take their animals out to the fields to feed on grass.


F) Mr Dittoh says the idea of converting online weather reports into speech came from the farmers themselves, after a workshop in the village of Guabuligah. “They came up with this,” he says. Mr Meteo takes the online weather forecast, converts it to a short recording in the appropriate language and makes it available on a basic phone. Farmers ring up to receive the information. The local language Dagbani is spoken by 1.2 million people but is not served by any online translation applications. The service was designed to be cheap and easy to run, says Mr Ditto. He plans to begin field tests this month, working with Tamale’s Savanna Agricultural Research Institute.


G) The spoken web could also help the one-in-five adults in Europe and the US with poor reading skills. But building the spoken web — web-to-voice and voice-to-web — isn’t straightforward. For software to understand pizza is served at Italian restaurants is easy. To cover multiple domains and to be able to have a conversation with users on every single topic is still a long way off.


H) So although many computer assistants can answer simple questions about the weather and play music for us, anything resembling a wide-ranging human conversation is decades away. Artificial intelligence just isn’t smart enough yet. Even turning your voice into text — automatic speech recognition — is one of the hardest problems to solve, as there are as many ways to pronounce things as there are people on the planet.


I) Siri has often been praised for its ability to interpret our casual language and deliver very specific and accurate result, sometimes even providing additional information. But it is still somewhat restricted, particularly when the language moved away from stiffer commands into more human interactions. In one example, the phrase “Send a text to Jason, Clint, Sam, and Lee saying we’re having dinner at Silver Cloud restaurant” was interpreted as sending a message to Jason only, containing the text “Clint Sam and Lee saying we’re having dinner at Silver Cloud restaurant”. It has also been noted that Siri lacks a proper editing function, as saying “Edit message to say: we’re at Silver Cloud restaurant and you should come find us” generates “Clint Sam and Lee saying we’re having dinner at Silver Cloud restaurant to say we’re at Silver Cloud restaurant and you should come find us”.


J) Using voice interaction feels far more intimate than surfing the net the old-fashioned way. This is intentional as the informal tone of the assistant helps create an emotional attachment. But if something speaks, it must also listen. Our phones are always near us and they are collecting data about us all the time. This has already raised privacy concerns. The American Civil Liberties Union has stated that digital assistants create a threat to privacy from hackers. Some people have other concerns. They worry assistants will one day be used to deliver advertising directly to us.


K) But digital voices need more personality to make them popular. Robots are not yet witty, Siri is boring. The benefits of using voice instead of tapping fingers obviously depend on the context. Doctors completing online forms about their patients by speech, for example, can dictate 150 words a minute, three times faster than typing on a keyboard. This enables them to spend less time on administration and more time with patients.


L) Last year, speech recognition company Nuance helped a doctors’ surgery in Dukinfield, near Manchester, set up a speech system for the practice’s six doctors. Now they can dictate notes on a patient’s health condition and treatment and a smart assistant automatically enters the information into the right fields on a web form. Previously, the doctors made voice recordings that were then transcribed by secretaries — a process that was costly and likely to cause delays. The new system means letters to patients now have more detail.


M) Using voice also makes sense when you’re doing other things with your hands. Think about when you’re cooking, and you just want to know what the next step in the recipe is. Your hands are covered with oil; you’re not going to get on the iPad, so it’s a lot more natural to talk. And speech obviously makes sense when you’re driving. In the US, 29\% of drivers admit they surf behind the wheel, according to insurance firm State Farm. This is up from 13\% in 2009. No wonder using mobile phones while driving causes more crashes a year than drink driving, says the US National Safety Council.


36. According to Francis Dittoh, their speech-based weather information service was meant to be inexpensive and easy to use.


37. Using voice instead of typing enables doctors to spend more time taking care of patients.


38. It is extremely difficult to convert voice into text because of different pronunciations.


39. African farmers unable to read often don’t have access to important information conveyed online.


40. Some phone users worry advertisers will take advantage of voice assistants to send ads directly to them.


41. The spoken web is helpful when one’s hands are occupied.


42. Some people believe online interaction would soon depend mainly on voice.


43. Setting up a spoken web is by no means an easy task.


44. Weather information is extremely important to farmers.


45. Some people are concerned about privacy because their phones are constantly collecting their personal information.


\subsection*{\textbf{Section C}}


\textbf{Directions: }\textit{There are 2 passages in this section. Each passage is followed by some questions or unfinished} \textit{statements. For each of them there are four choices marked A), B), C) and D). You should decide on the best} \textit{choice and mark the corresponding letter on }\textbf{\textit{Answer Sheet 2}}\textit{ with a single line through the centre.}


\subsection*{\textbf{Passage One}}


\textbf{Questions 46 to 50 are based on the following passage.}


The United States is facing a housing crisis: Affordable housing is inadequate, while luxury homes abound (充裕), and homelessness remains a persistent problem. Despite this, popular culture and the housing industry market happiness as living with both more space and more amenities (便利设施). Big houses are advertized as a reward for hard work and diligence, turning housing from a basic necessity into a luxury.


This is reflected in our homes. The average single-family home built in the United States before 1970 was less than 1,500 square feet in size. By 2016, the average size of a new, single-family home was 2,422 square feet. What’s more, homes built in the 2000s were more likely than earlier models to have more of all types of spaces: bedrooms, bathrooms, living rooms, dining rooms, recreation rooms and garages.


There are consequences of living big. As middle-class houses have grown larger, two things have happened. First, large houses take time to maintain, so cleaners and other low-wage service workers are required to keep these houses in order. Second, once-public spaces, where people from diverse backgrounds used to come together, have increasingly become privatized, leading to a reduction in the number of public facilities available to all, and a reduced quality of life for many. Take swimming pools. While in 1950, only 2,500 U.S. families owned pools, by 1999 this number was 4 million. At the same time, public municipal pools were often closed, leaving low-income people nowhere to swim.


The trend for bigger housing thus poses ethical questions. Should Americans accept a system in which the middle and upper classes enjoy a luxurious lifestyle, using the low-wage labor of others? Are we willing to accept a system in which an increase in amenities purchased by the affluent means a reduction in amenities for the poor?


I believe neither is acceptable. We must change the way we think: living well does not need to mean having more private spaces; instead, it could mean having more public spaces. A better goal than building bigger houses for some is to create more publicly accessible spaces and amenities for all.


46. What are big houses promoted to be in the United States?


\begin{itemize}[leftmargin=2em,itemsep=.2em,topsep=.4em]\item[A.] A luxury for the homeless.
\item[B.] A reward for industriousness.
\item[C.] An abundant source of comforts.
\item[D.] An absolute necessity for happiness.
\end{itemize}

47. What is one of the consequences of living big?


\begin{itemize}[leftmargin=2em,itemsep=.2em,topsep=.4em]\item[A.] Many Americans’ quality of life has become lower.
\item[B.] People from diverse backgrounds no longer socialize.
\item[C.] People no longer have access to public swimming pools.
\item[D.] Many Americans’ private life has been negatively affected.
\end{itemize}

48. What questions arise from living big?


\begin{itemize}[leftmargin=2em,itemsep=.2em,topsep=.4em]\item[A.] Questions related to moral principles.
\item[B.] Questions having to do with labor cost.
\item[C.] Questions about what lifestyle to promote.
\item[D.] Questions concerning housing development.
\end{itemize}

49. What kind of social system does the author think is unacceptable?


\begin{itemize}[leftmargin=2em,itemsep=.2em,topsep=.4em]\item[A.] One in which the wealthy exploit the low-wage laborers building their houses.
\item[B.] One in which the rich purchase amenities at an increasingly unjustifiable price.
\item[C.] One in which the upper classes deprive the lower classes of affordable housing.
\item[D.] One in which the affluent enjoy a more comfortable life at the expense of the poor.
\end{itemize}

50. What does the author advocate for people to live well?


\begin{itemize}[leftmargin=2em,itemsep=.2em,topsep=.4em]\item[A.] Finding ways to turn private spaces into public ones.
\item[B.] Building more houses affordable to those less affluent.
\item[C.] More public spaces created for everyone to enjoy.
\item[D.] All amenities made accessible to the rich and the poor alike.
\end{itemize}

\subsection*{\textbf{Passage Two}}


\textbf{Questions 51 to 55 are based on the following passage.}


Most of us in the entrepreneurial community are blessed — or cursed — with higher-than-average ambition. Ambitious people strongly desire accomplishments and are willing to take more risks and spend more effort to get them.


Overall, this is a positive quality, especially for people trying to build their own businesses. Apparently, if you’re more naturally driven to set goals, you are more likely to succeed.


Actually, this isn’t always the case. In fact, in some cases, extreme ambition may end up doing more harm than good.


One major side effect of excessive ambition is the tendency to focus too determinedly on one particular vision or end goal. This is problematic because it hinders your ability to adapt to new circumstances, which is vital if you want to be a successful entrepreneur. If a new competitor emerges to threaten your business, you may need to change direction, even if that means straying from your original vision. If you have too much ambition, you’ll find this hard, if not impossible.


Few people are successful when they try to build their first brand. Unfortunately, for the most ambitious entrepreneurs, a failure is seen as disastrous, and impossible to recover from.


It’s a clear departure from the intended plan toward the intended goal. For people with limited ambition, however, failure is viewed as something closer to reality. Remember, failure is inevitable, and every failure you survive is a learning experience.


Ambitious people tend to be more materialistically successful than their non-ambitious counterparts. However, they’re only slightly happier than their less-ambitious counterparts, and tend to live significantly shorter lives. This implies that even though ambitious people are more likely to achieve conventional “success,” such success means nothing for their health and happiness — and if you don’t have health and happiness. what else could possibly matter?


Clearly, some amount of ambition is good for your motivation. Without any ambition, you wouldn’t start your own business, set or achieve goals and get far in life. But an excess of ambition can also be dangerous, putting you at risk of burnout, stubbornness and even a shorter life.


51. What does the author think of most entrepreneurs?


\begin{itemize}[leftmargin=2em,itemsep=.2em,topsep=.4em]\item[A.] They are more willing to risk their own lives.
\item[B.] They are more ambitious than ordinary people.
\item[C.] They achieve greater nonconventional success.
\item[D.] They have more positive qualities than most of us.
\end{itemize}

52. What does the author imply by saying “this isn’t always the case” (Line 1, Para. 3)?


\begin{itemize}[leftmargin=2em,itemsep=.2em,topsep=.4em]\item[A.] Ambitious people may not have a greater chance of success.
\item[B.] Ambitious people may not have more positive qualities.
\item[C.] Entrepreneurs’ ambition does as much good as harm.
\item[D.] Entrepreneurs are more naturally driven to success.
\end{itemize}

53. What does the author say is of extreme importance for one to become a successful entrepreneur?


\begin{itemize}[leftmargin=2em,itemsep=.2em,topsep=.4em]\item[A.] Holding on to one’s original vision.
\item[B.] Being able to adapt to new situations.
\item[C.] Focusing determinedly on one particular goal.
\item[D.] Avoiding radical change in one’s career direction.
\end{itemize}

54. How do the most ambitious entrepreneurs regard failure in their endeavor?


\begin{itemize}[leftmargin=2em,itemsep=.2em,topsep=.4em]\item[A.] It will awaken them to reality.
\item[B.] It is a lesson they have to learn.
\item[C.] It means the end of their career.
\item[D.] It will result in a slow recovery.
\end{itemize}

55. What does the author advise us to do concerning ambition?


\begin{itemize}[leftmargin=2em,itemsep=.2em,topsep=.4em]\item[A.] Distinguish between conventional success and our life goal.
\item[B.] Follow the example of the most ambitious entrepreneurs.
\item[C.] Avoid taking unnecessary risks when starting a business.
\item[D.] Prioritize health and happiness over material success.
\end{itemize}

\subsection*{\textbf{Part IV Translation} \textbf{( 30 minutes)}}


\textbf{Directions: }\textit{For this part, you are allowed 30 minutes to translate a passage from Chinese into English. You} \textit{should write your answer on }\textbf{\textit{Answer Sheet 2}}\textit{.}


中国政府一直大力推行义务教育(compulsory education)，以使每个儿童都享有受教育的机会。 自1986 年《义务教育法》生效以来，经过不懈努力，实现了在全国推行义务教育的目标。如今， 在中国，儿童年满六周岁开始上小学，从小学到初中一共接受九年义务教育。从 2008 年秋季学期开始，义务教育阶段学生无需缴纳学费。随着一系列教育改革举措的实施，中国义务教育的质量也有显著提高。

\end{document}